% CP-VI-0036
% target_chapter: VI/26 — Finite-Type and Noetherian Morphisms
% source_unit_id: IMP-DL-D27460A55B-P005
% source: Downloads/theory-of-algebraic-geometry-5.tex L2846-L3060
% solution_provenance: SOURCE_EXPLICIT_SOLUTION

\subsection*{Problem 6: Examples violating the properties}

\textbf{Problem.}
For each of the properties defined above of schemes or morphisms, give an example of a scheme or morphism which violates that property. Give an example of a morphism which is locally of finite type but not of finite type.

\bigskip

\textbf{Solution.}

We list standard examples.

\subsection*{A scheme which is not reduced}

Let
\[
X=\operatorname{Spec} k[\varepsilon]/(\varepsilon^2).
\]
The ring
\[
k[\varepsilon]/(\varepsilon^2)
\]
has a nonzero nilpotent element, namely \(\varepsilon\), because
\[
\varepsilon\neq 0,
\qquad
\varepsilon^2=0.
\]
Therefore the ring is not reduced.

Hence
\[
\operatorname{Spec} k[\varepsilon]/(\varepsilon^2)
\]
is a scheme which is not reduced.

\subsection*{A scheme which is not irreducible}

Let
\[
X=\operatorname{Spec} k[x,y]/(xy).
\]
Geometrically, this is the union of the two coordinate axes in the affine plane.

Indeed, in \(\mathbb A_k^2=\operatorname{Spec} k[x,y]\), the equation
\[
xy=0
\]
means
\[
x=0
\qquad\text{or}\qquad
y=0.
\]
Thus
\[
\operatorname{Spec} k[x,y]/(xy)
\]
is the union of the closed subsets
\[
V(x)
\qquad\text{and}\qquad
V(y).
\]
Both are proper closed subsets. Therefore \(X\) is reducible.

Hence
\[
X=\operatorname{Spec} k[x,y]/(xy)
\]
is not irreducible.

\subsection*{A scheme which is not integral}

A scheme is integral if it is both reduced and irreducible.

Thus either a non-reduced scheme or a reducible scheme gives an example of a non-integral scheme.

For example,
\[
X=\operatorname{Spec} k[x,y]/(xy)
\]
is reduced but reducible. Therefore it is not integral.

Another example is
\[
X=\operatorname{Spec} k[\varepsilon]/(\varepsilon^2),
\]
which is not reduced. Hence it is not integral.

\subsection*{A morphism which is not locally of finite type}

Consider the morphism
\[
f:\operatorname{Spec} k[x_1,x_2,x_3,\dots]\longrightarrow \operatorname{Spec} k.
\]
This morphism corresponds to the inclusion
\[
k\longrightarrow k[x_1,x_2,x_3,\dots].
\]

The ring
\[
k[x_1,x_2,x_3,\dots]
\]
is not a finitely generated \(k\)-algebra, because it has infinitely many algebraically independent variables.

Since the target is affine,
\[
\operatorname{Spec} k,
\]
the definition of locally finite type would require the source to be covered by affine opens whose coordinate rings are finitely generated \(k\)-algebras.

But the whole source is
\[
\operatorname{Spec} k[x_1,x_2,x_3,\dots],
\]
and the algebra contains infinitely many independent variables. This is the standard example of a morphism which is not locally of finite type.

Therefore
\[
\operatorname{Spec} k[x_1,x_2,x_3,\dots]\longrightarrow \operatorname{Spec} k
\]
is not locally of finite type.

\subsection*{A morphism which is locally of finite type but not of finite type}

Let
\[
X=\coprod_{n\geq 1}\mathbb A_k^1
\]
be a countably infinite disjoint union of affine lines over \(k\). Consider the natural morphism
\[
f:X\longrightarrow \operatorname{Spec} k.
\]

Each component is
\[
\mathbb A_k^1=\operatorname{Spec} k[t],
\]
and \(k[t]\) is a finitely generated \(k\)-algebra. Therefore each component is of finite type over \(\operatorname{Spec} k\).

The inverse image of the only affine open of \(\operatorname{Spec} k\) is
\[
f^{-1}(\operatorname{Spec} k)=X.
\]
This is covered by the affine opens
\[
\mathbb A_k^1,\mathbb A_k^1,\mathbb A_k^1,\dots
\]
one for each component. On each such open, the coordinate ring is \(k[t]\), which is finitely generated over \(k\).

Therefore \(f\) is locally of finite type.

However, \(X\) cannot be covered by finitely many of these connected components. Indeed, a finite union of components would miss all the remaining components. Hence \(X\) is not quasi-compact.

Since a morphism of finite type is, in particular, quasi-compact, this morphism cannot be of finite type.

Therefore
\[
\coprod_{n\geq 1}\mathbb A_k^1
\longrightarrow
\operatorname{Spec} k
\]
is locally of finite type but not of finite type.

\subsection*{A scheme which is not a variety}

In the language of the remark, a variety over an algebraically closed field \(k\) is an integral scheme of finite type over \(\operatorname{Spec} k\).

Thus a scheme can fail to be a variety in several ways.

For example,
\[
\operatorname{Spec} k[\varepsilon]/(\varepsilon^2)
\]
is not a variety because it is not reduced.

Also,
\[
\operatorname{Spec} k[x,y]/(xy)
\]
is not a variety because it is reducible, hence not integral.

Finally,
\[
\operatorname{Spec} k[x_1,x_2,x_3,\dots]
\]
is not a variety because it is not of finite type over \(k\).

\subsection*{Summary of examples}

\[
\begin{array}{c|c}
	\textbf{Property violated} & \textbf{Example} \\
	\hline
	\text{Reduced} &
	\operatorname{Spec} k[\varepsilon]/(\varepsilon^2)
	\\[4pt]
	\text{Irreducible} &
	\operatorname{Spec} k[x,y]/(xy)
	\\[4pt]
	\text{Integral} &
	\operatorname{Spec} k[x,y]/(xy)
	\\[4pt]
	\text{Locally of finite type} &
	\operatorname{Spec} k[x_1,x_2,x_3,\dots]\to \operatorname{Spec} k
	\\[4pt]
	\text{Finite type} &
	\displaystyle \coprod_{n\geq 1}\mathbb A_k^1\to \operatorname{Spec} k
	\\[8pt]
	\text{Variety over }k &
	\operatorname{Spec} k[\varepsilon]/(\varepsilon^2)
\end{array}
\]
